\begin{abstract}
Operating Systems enforce logical isolation using abstractions such as processes, containers, and isolation technologies to protect a system from malicious or buggy code.  In this paper, we show new types of side channels through the file system that break this logical isolation.  The file system plays a critical role in the operating system, managing all I/O activities between the application layer and the physical storage device. 
%Prior research has identified I/O side-channel vulnerabilities at both the physical and system levels. Early studies demonstrated that I/O pins can serve as leakage vectors for physical side channels, while others showed that contention on I/O can also lead to leakage. However, these attacks are often constrained by the need for physical access or specific application designs. 
We observe that the file system implementation is shared, leading to timing leakage when using common I/O system calls. Specifically, we found that modern operating systems take advantage of any flush operation (which saves cached blocks in memory to the SSD or disk) to flush all of the I/O buffers, even those used by other isolation domains.   Thus, by measuring the delay of \textit{syncfs}, the attacker can infer the I/O behavior of victim programs.  We then demonstrate a \textit{syncfs} covert channel attack on multiple file systems, including both Linux native file systems and the Windows file system, achieving a maximum bandwidth of 5 Kbps with an error rate of 0.15\% on Linux and 7.6 Kbps with an error rate of 1.9\% on Windows. 
In addition, we construct three side-channel attacks targeting both Linux and Android devices. On Linux devices, we implement a website fingerprinting attack and a video fingerprinting attack by tracking the write patterns of temporary buffering files. On Android devices, we design an application fingerprinting attack that leaks application write patterns during boot-up.  The attacks achieve over 90\% F1 score, precision, and recall. Finally, we demonstrate that these attacks can be exploited across containers implementing a container detection technique and a cross-container covert channel attack.


% Additionally, we construct cross-container attacks with \textit{syncfs} including a container detection technique and a cross-container covert channel attack.
% Besides that, we also build three side channel attacks on both Linux and Android devices. On Linux devices, we construct a website fingerprinting attack and a video fingerprinting attack by leaking the write pattern of temporary files. On Android devices, we construct an application fingerprinting attack by leaking the application write pattern during boot-up. The performance of these attacks is above 90\% for F1, precision, and recall.

\end{abstract}